  \item \field{Innovation Concept and Background}

    \textbf{The problem.} A person cannot hold a job if they cannot reliably manage
    the ordinary structure of a day --- taking medication on time, remembering an
    appointment, following the steps of a task in order. For adults living with
    cognitive and memory disability, this everyday structure is precisely what is
    impaired. Yet in Thailand this population is largely locked out of work. Of
    roughly 855{,}000 working-age people with disability, only about
    \textbf{36.3\%} are employed, and merely \textbf{8\%} participate in the
    competitive labour market~[3, 4]. Cognitive-type disability --- intellectual
    disability (about 145{,}000 registered) and mental or behavioural disability
    (about 167{,}000 registered) --- makes up a large and especially under-served
    share of this group~[2].

    This exclusion is caused by the disability, not incidental to it. Across
    traumatic brain injury, post-stroke impairment, and intellectual disability,
    deficits in memory, attention, and executive function are the strongest
    predictors of whether a person can return to and keep a job~[5, 7, 8] --- after
    traumatic brain injury, unemployment can reach 60\% two years post-injury~[8].
    The barrier to earning, for this group, is a barrier to remembering and organising.

    \textbf{Target group.} Working-age adults (15--59) with cognitive or memory
    disability: intellectual disability, traumatic brain injury, post-stroke
    cognitive impairment, and early-onset dementia.

    \textbf{Why current solutions fall short.} Thailand's disability allowance
    (800 baht/month, roughly 29\% of the poverty line) is income support, not a
    route to work~[1]. The employment quota law is weakly enforced, and employers
    frequently pay fines rather than hire~[6]. The assistive technology that does
    target cognition is almost entirely phone- or tablet-based --- which fails the
    very users it is meant to help: it requires them to remember to open an app, to
    carry and charge a device, and to navigate menus. Abandonment is high when
    setup is complex or support is absent~[14, 17, 18]. The memory aid should not
    itself demand memory to use.

    \textbf{Why an ambient smart mirror, and why now.} A mirror is furniture people
    already pass every day: it needs no activation ritual, nothing to carry, and does
    not look like a medical device. It can offer support at the exact place and moment
    it is needed --- the morning medication routine at the bathroom mirror, an
    orientation to the day, a step-by-step task prompt. This is not speculative: a
    smart-mirror cognitive-training system has already been built and shown to improve
    cognitive scores, recall, and orientation in adults with cognitive disability~[9],
    smart mirrors for independent living have been demonstrated~[10], and randomised
    evidence shows prompting technology reduces the assistance people with acquired
    brain injury need for daily activities~[11]. The enabling technologies ---
    reliable on-device voice, vision, and low-cost displays --- have only recently
    matured enough to make an affordable, private, always-on mirror practical.

    \textbf{From Learning to Earning.} Our thesis is a causal chain. Reliable daily
    routines make a person independent; independence makes them able to learn and
    re-learn work tasks; the ability to learn tasks makes them employable. The mirror
    is an ambient coach across that whole chain --- it steadies the daily foundation
    and then carries the same step-by-step prompting into work skills. Assistive
    technology has been shown in systematic review to improve employment outcomes for
    people with cognitive disability~[14]. We aim to turn daily independence into a
    genuine path to earning.
